\chapter{CONTROL DESIGN}
\label{ch:control}

The TWIP is unstable and must be actively controlled. Linear control is
effective near the upright equilibrium when the parameters are accurately known,
so a linear quadratic regulator is used as the nominal control. An additive
signal is then formulated to compensate for the unmatched uncertainties
appearing in the two acceleration equations of \eqref{eq:plantform}.

\section{LQR CONTROL DESIGN}
\label{sec:lqr}

Because the controller is implemented digitally, the discrete-time infinite
horizon formulation is used. The performance index
\begin{equation}\label{eq:lqrJ}
J=\sum_{k=0}^{\infty}\bigl(x_k^TQ_{\mathrm{L}}x_k+u_k^TR_{\mathrm{L}}u_k\bigr)
\end{equation}
is minimized by solving the discrete algebraic Riccati equation
\cite{simon2006optimal}
\begin{equation}\label{eq:dare}
P=Q_{\mathrm{L}}+F^T\Bigl(P-PG\bigl(R_{\mathrm{L}}+G^TPG\bigr)^{-1}G^TP\Bigr)F,
\end{equation}
with $F$ and $G$ from the discretized model of Section~\ref{sec:discretization}
and $Q_{\mathrm{L}}\succeq0$, $R_{\mathrm{L}}\succ0$ the weighting matrices.
\rev{Note that \eqref{eq:dare} carries $F$ on both sides; revision~7 wrote the
trailing factor as $A$, which is a typographical carryover from the
continuous-time notation and is corrected here.} The control law is
\begin{equation}\label{eq:lqrlaw}
u_{\mathrm{nom},k}=-K_{\mathrm{LQR}}\bigl(x_k-x_{\mathrm{des},k}\bigr),
\qquad
K_{\mathrm{LQR}}=\bigl(R_{\mathrm{L}}+G^TPG\bigr)^{-1}G^TPF .
\end{equation}

\section{UNMATCHED UNCERTAINTY AND THE VIRTUAL CONTROL}

When the uncertainty enters the same channel as the control, the system is said
to have matched uncertainty, and an estimate of that uncertainty can be used
directly to cancel it. The TWIP does not have this structure. In
\eqref{eq:plantform} the scalar input $u$ enters both $\dot x_2$ and $\dot x_4$
through $B_1$ and $B_2$, while two independent uncertainties $d_1$ and $d_2$
appear in the same two equations. No choice of $u$ cancels both.

The device used here, following Huang \cite{huang2002lyapunov} and related in
spirit to the partial feedback linearization of Pathak et al.\
\cite{pathak2005twip}, is to treat
the tilt state $x_3$ as a virtual control for the position subsystem
$(x_1,x_2)$, and then to design the actual control $u$ so that $x_3$ tracks the
value the position subsystem requires. The construction below is a four-step
backstepping procedure with command filtering.

\rev{\subsection{Three requirements on the construction}

Three requirements shape the design and distinguish it from a two-step version.

\begin{enumerate}[label=(R\arabic*),leftmargin=*]

\item \emph{The coupling term must be retained.} Substituting the desired value
of $x_3$ into $\dot x_2$ is valid only if $x_3$ equals that value. It does not;
the difference is precisely the error coordinate $z_3$ introduced at the next
step. Carrying $x_3=z_3+x_3^c$ through leaves an $A_2z_3$ term in $\dot z_2$
that must be cancelled explicitly. Dropping it removes a term of the same order
as the ones being cancelled.

\item \emph{Every error coordinate needs damping.} If $z_1$ and $z_3$ appear in
the Lyapunov function but not in its derivative, uniform ultimate boundedness
cannot be concluded for them: the standard argument requires the derivative to
be negative outside a compact set in the full error space, and a coordinate
absent from the derivative fails this in its own direction. Concretely,
$\dot z_1=z_2$ is a pure integrator, so a merely bounded $z_2$ permits $z_1$ to
ramp. Damping is assigned to all four coordinates below.

\item \emph{Each network's target must be a function of measurable signals
alone.} If the function $F_1$ that network one is required to approximate
contains the control $u$, and $u$ contains the output of network two, and the
target $F_2$ of network two contains derivatives of the output of network one,
then neither target is a fixed continuous function of its network input and the
universal approximation hypothesis does not apply. Command filtering removes
this dependence: the filter outputs are states, so the required signals are
available at the current instant without differentiating any adaptive law.

\end{enumerate}}

\section{COMMAND FILTERS}
\label{sec:cmdfilter}

For each virtual control $\alpha_i$, $i=1,2,3$, define a second-order command
filter with bandwidth $\varpi_i>0$ and damping $\varsigma_i\in(0,1]$:
\begin{equation}\label{eq:cmdfilt}
\dot q_{i,1}=q_{i,2},\qquad
\dot q_{i,2}=-2\varsigma_i\varpi_i q_{i,2}-\varpi_i^2\bigl(q_{i,1}-\alpha_i\bigr),
\end{equation}
with filtered command $x_{i+1}^c:=q_{i,1}$ and its derivative
$\dot x_{i+1}^c:=q_{i,2}$, both available as filter states. Magnitude, rate, and
bandwidth limits may be imposed on \eqref{eq:cmdfilt} to respect actuator
constraints. Define the filter errors
\begin{equation}\label{eq:chi}
\chi_i:=x_{i+1}^c-\alpha_i .
\end{equation}

\begin{assumption}[Filter error]\label{as:filt}
On the compact operating set, the virtual controls $\alpha_i$ and their first
derivatives are bounded, and for each $i$ there exists $\bar\chi_i>0$ with
$\abs{\chi_i}\le\bar\chi_i$ and $\bar\chi_i=O(1/\varpi_i)$.
\end{assumption}

Assumption~\ref{as:filt} is the standard command-filtered backstepping result of
Farrell et al.\ \cite{farrell2009command} and Swaroop et al.\
\cite{swaroop2000dsc}: the filter error can be made
arbitrarily small by raising the bandwidth, subject to the usual separation from
the unmodeled dynamics and the sample rate.

\section{COMMAND-FILTERED NEURAL BACKSTEPPING}
\label{sec:cffb}

Let $x_{1d}(t)$ be the position reference, assumed twice continuously
differentiable with bounded derivatives. Write $u=u_{\mathrm{nom}}+u_e$ as in
Section~\ref{sec:lqr}.

\paragraph{Step 1.} Define $z_1:=x_1-x_{1d}$ and the virtual control
\begin{equation}\label{eq:alpha1}
\alpha_1:=\dot x_{1d}-c_1z_1 ,
\end{equation}
filtered to $x_2^c$. With $z_2:=x_2-x_2^c$,
\begin{equation}\label{eq:z1dot}
\dot z_1=x_2-\dot x_{1d}=z_2+\chi_1-c_1z_1 .
\end{equation}

\paragraph{Step 2.} With $x_3=z_3+x_3^c=z_3+\chi_2+\alpha_2$,
\begin{equation}\label{eq:z2raw}
\dot z_2=A_1x_2+A_2\alpha_2+A_2\bigl(z_3+\chi_2\bigr)+B_1u+d_1-\dot x_2^c .
\end{equation}
Define the first approximation target
\begin{equation}\label{eq:F1}
F_1:=A_1x_2+d_1(\mathbf{x})-\dot x_2^c ,
\end{equation}
which depends only on the measured state and on the filter state $\dot x_2^c$,
and in particular \emph{not} on the control. Choose
\begin{equation}\label{eq:alpha2}
\alpha_2:=A_2^{-1}\Bigl[-\hat F_1-B_1u-c_2z_2-z_1\Bigr].
\end{equation}
Substituting into \eqref{eq:z2raw},
\begin{equation}\label{eq:z2dot}
\dot z_2=\bigl(F_1-\hat F_1\bigr)+A_2z_3+A_2\chi_2-c_2z_2-z_1 .
\end{equation}
\rev{The term $A_2z_3$ is retained; it is cancelled at Step~3.}

\paragraph{Step 3.} With $z_3:=x_3-x_3^c$ and $x_4=z_4+x_4^c=z_4+\chi_3+\alpha_3$,
choose
\begin{equation}\label{eq:alpha3}
\alpha_3:=\dot x_3^c-c_3z_3-A_2z_2 ,
\end{equation}
filtered to $x_4^c$, giving
\begin{equation}\label{eq:z3dot}
\dot z_3=z_4+\chi_3-c_3z_3-A_2z_2 .
\end{equation}
The term $-A_2z_2$ in \eqref{eq:alpha3} is what cancels $+A_2z_3$ in
\eqref{eq:z2dot} when the two are combined in the Lyapunov derivative.

\paragraph{Step 4.} With $z_4:=x_4-x_4^c$, define the second approximation target
\begin{equation}\label{eq:F2}
F_2:=A_3x_2+A_4x_3+d_2(\mathbf{x})-\dot x_4^c ,
\end{equation}
again free of the control, and choose the extra control
\begin{equation}\label{eq:ue}
u_e:=B_2^{-1}\Bigl[-\hat F_2-B_2u_{\mathrm{nom}}-c_4z_4-z_3\Bigr],
\end{equation}
so that
\begin{equation}\label{eq:z4dot}
\dot z_4=\bigl(F_2-\hat F_2\bigr)-c_4z_4-z_3 .
\end{equation}

\rev{\begin{remark}[Execution order, and why there is no algebraic loop]
\label{rem:order}
At each sample the following order is used:
(1) read the state and the filter states $q_{i,\cdot}$;
(2) form $z_1,\dots,z_4$ from the state and the filter states;
(3) evaluate $\hat F_1,\hat F_2$;
(4) compute $u_{\mathrm{nom}}$ from \eqref{eq:lqrlaw} and $u_e$ from
\eqref{eq:ue}, hence $u$ (note that substituting \eqref{eq:ue} gives
$B_2u=-\hat F_2-c_4z_4-z_3$, so $u_{\mathrm{nom}}$ cancels and the LQR gain
does not affect the applied control);
(5) compute $\alpha_1,\alpha_2,\alpha_3$, with $\alpha_2$ using the $u$ just
computed;
(6) propagate \eqref{eq:cmdfilt} one step.
Although $\alpha_2$ depends on $u$ and $u$ depends on $z_3$, $z_3$ is formed
from the filter state $x_3^c$ rather than from $\alpha_2$ itself. The
dependence is therefore through a dynamic element and is resolved by the
ordering above. Without command filtering, $x_3^c$ would be replaced by
$\alpha_2$ and steps (4) and (5) would be mutually defining.
\end{remark}}

\section{NEURAL NETWORK APPROXIMATION}

A single-layer network was found inadequate for $F_1$ and $F_2$, so two-layer
networks are used. For $i=1,2$,
\begin{equation}\label{eq:nn2layer}
F_i=W_{i2}^T\sigma_{i2}\bigl(W_{i1}^T\sigma_{i1}(V_i^TP_i)\bigr)+\varepsilon_i,
\qquad \abs{\varepsilon_i}\le\varepsilon_{im},
\end{equation}
with $V_i$ a fixed randomly selected input layer, $P_i$ the network input
vector, and $\sigma_{i1},\sigma_{i2}$ bounded activation functions. The
estimates are
\begin{equation}\label{eq:nnhat}
\hat F_i=\What_{i2}^T\hat\sigma_{i2}
\bigl(\What_{i1}^T\hat\sigma_{i1}(V_i^TP_i)\bigr).
\end{equation}
Network inputs are taken as
$P_1=\bigl[x_2,\ x_3,\ \dot x_2^c,\ 1\bigr]^T$ and
$P_2=\bigl[x_2,\ x_3,\ x_4,\ \dot x_4^c,\ 1\bigr]^T$; by construction neither
contains a control signal or an adaptive parameter.

Writing $\Wt_{ij}:=W_{ij}-\What_{ij}$ and expanding the hidden-layer activation
about $\What_{i1}^T\hat\sigma_{i1}$ in the standard manner
\cite{lewis1999neural},
\begin{equation}\label{eq:nnerr}
F_i-\hat F_i
=\Wt_{i2}^T\Bigl(\hat\sigma_{i2}-\hat\sigma_{i2}'\What_{i1}^T\hat\sigma_{i1}\Bigr)
+\What_{i2}^T\hat\sigma_{i2}'\,\Wt_{i1}^T\hat\sigma_{i1}
+w_i ,
\end{equation}
where $\hat\sigma_{i2}'$ is the Jacobian of $\sigma_{i2}$ evaluated at the
estimate and $w_i$ collects $\varepsilon_i$ with the higher-order Taylor
residual. \rev{That residual contains the weight errors, so $w_i$ is not bounded
by a constant. Because the activations and their derivatives are bounded and the
ideal weights satisfy $\norm{W_{ij}}_F\le W_{ijm}$, it satisfies, on the compact
operating set,
\begin{equation}\label{eq:wbound}
\abs{w_i}\le w_{i0}+w_{i1}\norm{\Wt_i}_F,
\qquad \norm{\Wt_i}_F^2:=\norm{\Wt_{i1}}_F^2+\norm{\Wt_{i2}}_F^2,
\end{equation}
with constants $w_{i0},w_{i1}\ge0$ \cite{lewis1999neural}.}

\rev{Equation \eqref{eq:nnerr} is the form that makes both layers' weight errors
appear explicitly and linearly. Bounding the hidden-layer activation error by a
constant instead, as in revision~7, leaves the first-layer weight error
unmatched in the Lyapunov derivative, so that the first-layer update law is
unjustified by the analysis and the network is effectively single-layer as far
as the proof is concerned.}

Let $\varrho_1:=z_2$ and $\varrho_2:=z_4$ denote the error coordinate driving
network $i$. The weight update laws, each with $\sigma$-modification, are
\begin{align}
\dot{\What}_{i2}&=\gamma_{i2}\Bigl[
\bigl(\hat\sigma_{i2}-\hat\sigma_{i2}'\What_{i1}^T\hat\sigma_{i1}\bigr)\varrho_i^T
-\kappa_{i2}\What_{i2}\Bigr],
\label{eq:upd2}\\[4pt]
\dot{\What}_{i1}&=\gamma_{i1}\Bigl[
\hat\sigma_{i1}\,\varrho_i^T\,\What_{i2}^T\hat\sigma_{i2}'
-\kappa_{i1}\What_{i1}\Bigr],
\label{eq:upd1}
\end{align}
with $\gamma_{ij}>0$ the learning rates and $\kappa_{ij}>0$ the leakage gains of
the $\sigma$-modification \cite{ioannou1983adaptive,ioannou1984sigma}.
\rev{Both laws now carry the same sign convention; the mismatch between the two
layers in revision~7 is removed.}

\section{STABILITY ANALYSIS}

\begin{theorem}[Closed-loop boundedness]\label{thm:control}
Consider \eqref{eq:plantform} with the control \eqref{eq:lqrlaw},
\eqref{eq:ue}, the virtual controls \eqref{eq:alpha1}, \eqref{eq:alpha2},
\eqref{eq:alpha3}, the command filters \eqref{eq:cmdfilt}, and the update laws
\eqref{eq:upd2}--\eqref{eq:upd1}. Let Assumption~\ref{as:filt} hold, let the
ideal weights satisfy $\norm{W_{ij}}_F\le W_{ijm}$, and let the damping gains
satisfy
\begin{equation}\label{eq:gaincond}
c_1>\tfrac{1}{2},\qquad
c_2>1,\qquad
c_3>\tfrac{1}{2},\qquad
c_4>\tfrac{1}{2},\qquad
\kappa_{ij}>2w_{i1}^2 ,
\end{equation}
\rev{the first four after the cross terms are absorbed as below, the last so
that the leakage dominates the weight-error part of \eqref{eq:wbound}.} Then all tracking errors
$z_1,\dots,z_4$ and all weight estimation errors $\Wt_{ij}$ are uniformly
ultimately bounded.
\end{theorem}

\begin{proof}
Take
\begin{equation}\label{eq:Vctrl}
\mathcal{V}=\frac{1}{2}\sum_{i=1}^{4}z_i^Tz_i
+\frac{1}{2}\sum_{i=1}^{2}\sum_{j=1}^{2}
\frac{1}{\gamma_{ij}}\Tr\bigl(\Wt_{ij}^T\Wt_{ij}\bigr).
\end{equation}
Differentiating and substituting
\eqref{eq:z1dot}, \eqref{eq:z2dot}, \eqref{eq:z3dot}, \eqref{eq:z4dot},
\[
\begin{aligned}
\dot{\mathcal{V}}=&-c_1z_1^2-c_2z_2^2-c_3z_3^2-c_4z_4^2\\
&+\underbrace{z_1z_2-z_2z_1}_{=0}
+\underbrace{A_2z_2z_3-A_2z_3z_2}_{=0}
+\underbrace{z_3z_4-z_4z_3}_{=0}\\
&+z_1\chi_1+A_2z_2\chi_2+z_3\chi_3
+z_2\bigl(F_1-\hat F_1\bigr)+z_4\bigl(F_2-\hat F_2\bigr)
-\sum_{i,j}\frac{1}{\gamma_{ij}}\Tr\bigl(\Wt_{ij}^T\dot{\What}_{ij}\bigr).
\end{aligned}
\]
The three bracketed groups vanish identically: the first because $\alpha_1$
contributes $-z_1$ to $\dot z_2$ via \eqref{eq:alpha2}, the second because of
the $-A_2z_2$ term in \eqref{eq:alpha3}, and the third because of the $-z_3$
term in \eqref{eq:ue}.

Substituting \eqref{eq:nnerr} for $F_i-\hat F_i$, and writing
$\Lambda_i:=\hat\sigma_{i2}-\hat\sigma_{i2}'\What_{i1}^T\hat\sigma_{i1}$ for the
effective second-layer regressor, the approximation terms become
\begin{equation}\label{eq:approxterms}
\varrho_i^T\bigl(F_i-\hat F_i\bigr)
=\varrho_i^T\Wt_{i2}^T\Lambda_i
+\varrho_i^T\What_{i2}^T\hat\sigma_{i2}'\Wt_{i1}^T\hat\sigma_{i1}
+\varrho_i^Tw_i .
\end{equation}
Each of the first two is a scalar and may be written as a trace,
\begin{align}
\varrho_i^T\Wt_{i2}^T\Lambda_i
&=\Tr\bigl(\Wt_{i2}^T\Lambda_i\varrho_i^T\bigr),\label{eq:tr1}\\[2pt]
\varrho_i^T\What_{i2}^T\hat\sigma_{i2}'\Wt_{i1}^T\hat\sigma_{i1}
&=\Tr\bigl(\Wt_{i1}^T\hat\sigma_{i1}\varrho_i^T\What_{i2}^T\hat\sigma_{i2}'\bigr).
\label{eq:tr2}
\end{align}
The weight error contributions to $\dot{\mathcal{V}}$ follow from
$\dot{\Wt}_{ij}=-\dot{\What}_{ij}$ and \eqref{eq:upd2}--\eqref{eq:upd1},
\begin{align}
\frac{1}{\gamma_{i2}}\Tr\bigl(\Wt_{i2}^T\dot{\Wt}_{i2}\bigr)
&=-\Tr\bigl(\Wt_{i2}^T\Lambda_i\varrho_i^T\bigr)
+\kappa_{i2}\Tr\bigl(\Wt_{i2}^T\What_{i2}\bigr),\label{eq:wdot2}\\[2pt]
\frac{1}{\gamma_{i1}}\Tr\bigl(\Wt_{i1}^T\dot{\Wt}_{i1}\bigr)
&=-\Tr\bigl(\Wt_{i1}^T\hat\sigma_{i1}\varrho_i^T\What_{i2}^T\hat\sigma_{i2}'\bigr)
+\kappa_{i1}\Tr\bigl(\Wt_{i1}^T\What_{i1}\bigr).\label{eq:wdot1}
\end{align}
Comparing \eqref{eq:tr1} with \eqref{eq:wdot2} and \eqref{eq:tr2} with
\eqref{eq:wdot1}, every \rev{explicit} term containing both a weight error and an error
coordinate appears exactly twice with opposite sign and cancels. \rev{(The
residual $w_i$ still depends on the weight errors through \eqref{eq:wbound};
it is handled with the leakage below.) This
pairwise cancellation is the purpose of the specific regressors chosen in
\eqref{eq:upd2}--\eqref{eq:upd1}, and it is available only because
\eqref{eq:nnerr} keeps both layers' weight errors explicit and linear. Bounding
the hidden-layer activation error by a constant leaves \eqref{eq:tr2} with no
counterpart, and the first-layer update law then contributes nothing that the
analysis can use.}

What survives is the leakage,
\[
\sum_{i,j}\kappa_{ij}\Tr\bigl(\Wt_{ij}^T\What_{ij}\bigr)
=\sum_{i,j}\kappa_{ij}\Tr\bigl(\Wt_{ij}^T(W_{ij}-\Wt_{ij})\bigr)
\le\sum_{i,j}\Bigl[-\frac{\kappa_{ij}}{2}\norm{\Wt_{ij}}_F^2
+\frac{\kappa_{ij}}{2}W_{ijm}^2\Bigr],
\]
together with the filter-error and residual cross terms. Hence
\[
\dot{\mathcal{V}}\le
-\sum_{i=1}^{4}c_iz_i^2
-\sum_{i,j}\frac{\kappa_{ij}}{2}\norm{\Wt_{ij}}_F^2
+z_1\chi_1+A_2z_2\chi_2+z_3\chi_3+z_2w_1+z_4w_2
+\sum_{i,j}\frac{\kappa_{ij}}{2}W_{ijm}^2 .
\]
Applying $ab\le\tfrac12a^2+\tfrac12b^2$ to each cross term, \rev{and for the
residuals
$\varrho_iw_i\le\tfrac12\varrho_i^2+\tfrac12\bigl(w_{i0}+w_{i1}\norm{\Wt_i}_F\bigr)^2
\le\tfrac12\varrho_i^2+w_{i0}^2+w_{i1}^2\norm{\Wt_i}_F^2$,}
\begin{equation}\label{eq:Vdotfinal}
\dot{\mathcal{V}}\le
-\Bigl(c_1-\tfrac12\Bigr)z_1^2
-\Bigl(c_2-1\Bigr)z_2^2
-\Bigl(c_3-\tfrac12\Bigr)z_3^2
-\Bigl(c_4-\tfrac12\Bigr)z_4^2
-\sum_{i,j}\Bigl(\frac{\kappa_{ij}}{2}-w_{i1}^2\Bigr)\norm{\Wt_{ij}}_F^2
+\mathcal{C},
\end{equation}
where
\[
\mathcal{C}:=\tfrac12\bar\chi_1^2+\tfrac12A_2^2\bar\chi_2^2
+\tfrac12\bar\chi_3^2+w_{10}^2+w_{20}^2
+\sum_{i,j}\frac{\kappa_{ij}}{2}W_{ijm}^2 .
\]
With $c_2>1$, $c_1,c_3,c_4>\tfrac12$, \rev{and $\kappa_{ij}>2w_{i1}^2$}, every coefficient in
\eqref{eq:Vdotfinal} is negative and $\dot{\mathcal{V}}<0$ outside a compact set
in $(z_1,z_2,z_3,z_4,\Wt_{11},\Wt_{12},\Wt_{21},\Wt_{22})$. Uniform ultimate
boundedness follows by the standard Lyapunov argument, with ultimate bound
proportional to $\sqrt{\mathcal{C}}$.
\end{proof}

\rev{\begin{remark}[What the bound depends on]
The residual $\mathcal{C}$ has three distinct contributions, and they can be
reduced independently. The filter terms $\bar\chi_i^2$ shrink as the filter
bandwidths $\varpi_i$ are raised, limited by the sample rate and by the
unmodeled structural modes of the chassis. The network residuals
\rev{$w_{i0}^2$} shrink with the richness of the basis. The leakage terms $\kappa_{ij}W_{ijm}^2$
shrink with $\kappa_{ij}$, at the cost of weaker protection against weight
drift. This decomposition gives the tuning order recommended in
Appendix~\ref{app:validation}.
\end{remark}}

\rev{\begin{remark}[Realizability]\label{rem:realizability}
Theorem~\ref{thm:control} is conditional on Assumption~\ref{as:filt}, and for
the TWIP that assumption does not follow from the construction. The control
enters $\dot x_2$ as well as $\dot x_4$, so $\alpha_2$ contains $B_1u$, and $u$
depends on the filter states. The virtual controls and their derivatives
therefore depend on the filter errors that the assumption requires to be small.
With $x_1$ as the tracked output, the recursion inverts the plant from $u$ to
$x_1$, which has a right-half-plane zero. Section~\ref{sec:cfinstab} shows that
the implemented closed loop has an unstable eigenvalue for every gain set
tested, approaching that zero as the filter bandwidths are raised. Section~\ref{sec:flatcf} gives the corrected design.
\end{remark}}

\rev{\begin{remark}[Relation to the two-step design]\label{rem:twostep}
If the $A_2z_3$ term is dropped from \eqref{eq:z2dot} and no damping is assigned
to $z_1$ and $z_3$, the surviving negative terms in \eqref{eq:Vdotfinal} are
those in $z_2$ and $z_4$ only. The resulting statement bounds $z_2$ and $z_4$
but not $z_1$ and $z_3$, and since $\dot z_1=z_2$ is a pure integrator, a
bounded but nonvanishing $z_2$ permits unbounded $z_1$. In practice the position
loop is held by the LQR term, but because $u_{\mathrm{nom}}$ is absorbed into
the approximation targets in that formulation, the analysis does not credit it.
The design above makes the position damping explicit.
\end{remark}}

\rev{\section{CORRECTED DESIGN: FLAT-OUTPUT COORDINATES}
\label{sec:flatcf}

Remark~\ref{rem:realizability} traces the failure of Section~\ref{sec:cffb} to
the choice of $x_1$ as the tracked output and to the control appearing in
$\dot x_2$. Both are removed by a change of coordinates that leaves the rest of
the construction intact. The same four steps, command filters, damping
assignment, and two networks are used.

\subsection{Coordinates}

Let $b:=B_1/B_2$ and define
\begin{equation}\label{eq:flatcoords}
y:=x_1-b\,x_3,\qquad \eta:=\dot y=x_2-b\,x_4 .
\end{equation}
From \eqref{eq:plantform},
\begin{equation}\label{eq:etadot}
\dot\eta=(A_1-bA_3)\,x_2+(A_2-bA_4)\,x_3+D_1(\mathbf{x}),
\qquad D_1:=d_1-b\,d_2 ,
\end{equation}
which contains no control. For the TWIP the coefficient of $x_2$ vanishes
identically, as do the tilt-rate terms of the corrected model
(Section~\ref{sec:simmodel}). The back-EMF and friction act through the same
generalized-force direction as the voltage, so they cancel together with the
input. \rev{Physically, $(r\beta+M_pl)\eta$ is the angular momentum of the robot
about the wheel contact point, to first order (Remark~\ref{rem:momentum}). The
motor torque, being internal, cannot change it; only gravity does, and
$(r\beta+M_pl)a_2x_3=-M_pglx_3$ is the gravity moment about that point.} Hence
\begin{equation}\label{eq:etadot2}
\dot\eta=a_2x_3+D_1(\mathbf{x}),\qquad a_2:=A_2-bA_4 .
\end{equation}
The output $y$ has relative degree four and no zero dynamics; it is the flat
output of the linearization. Because $b$ is a ratio of input gains, the motor
constants cancel from both $b$ and $a_2$, which depend only on the masses,
geometry, and inertia\rev{:
\begin{equation}\label{eq:ba2}
b=\frac{I_p+M_pl^2+M_plr}{r\beta+M_pl},
\qquad
a_2=-\frac{M_pgl}{r\beta+M_pl}.
\end{equation}
For the parameters of Chapter~\ref{ch:platform}, with $I_p$ about the center of
gravity, $b=\SI{0.159}{m}$ and $a_2=-5.62~\mathrm{m\,s^{-2}}$.} The design
requires $a_2\neq0$, which holds for any $l>0$.

Note the sign. With the tilt held, $\dot x_4=0$ requires $B_2u=-A_4x_3$, and the
translational acceleration is then $a_2x_3$, not $A_2x_3$. The quasi-static
effect of tilt on acceleration is opposite in sign to $A_2$ \rev{and 1.6 times
larger in magnitude ($a_2=-5.62$ against $A_2=3.43$).}
The design of Section~\ref{sec:cffb} uses $A_2$ as the virtual-control gain and
relies on the $B_1u$ term in $\alpha_2$ to make up the difference
instantaneously. That is where the plant inversion occurs.

\subsection{Steps}

With the command filters of Section~\ref{sec:cmdfilter}, and $y_d$ the reference
for $y$ (for a constant velocity command $y_d=x_{1d}-b\,x_3^{ss}$),
\begin{align}
z_1&:=y-y_d, & \alpha_1&:=\dot y_d-c_1z_1,\label{eq:fz1}\\
z_2&:=\eta-\eta^c, & \alpha_2&:=a_2^{-1}\bigl[-\hat F_1-c_2z_2-z_1\bigr],\label{eq:fz2}\\
z_3&:=x_3-x_3^c, & \alpha_3&:=\dot x_3^c-c_3z_3-\tfrac{a_2}{w}\,z_2,\label{eq:fz3}\\
z_4&:=x_4-x_4^c, & u&:=B_2^{-1}\bigl[-\hat F_2-c_4z_4-z_3\bigr],\label{eq:fz4}
\end{align}
with $F_1:=D_1-\dot\eta^c$ and $F_2$ as in \eqref{eq:F2}. \rev{Accordingly the
input of network one is $P_1=\bigl[x_2,\ x_3,\ \dot\eta^c,\ 1\bigr]^T$, with the
filter state $\dot\eta^c$ in place of $\dot x_2^c$.} The control now enters
only at the last step, so $\alpha_2$ contains no control and the ordering of
Remark~\ref{rem:order} has no algebraic dependence at all. As before,
$u_{\mathrm{nom}}$ cancels if $u$ is written as $u_{\mathrm{nom}}+u_e$.

\subsection{Weighted Lyapunov function}

The constant $w>0$ in \eqref{eq:fz3} weights the attitude coordinates in
\begin{equation}\label{eq:Vflat}
\mathcal{V}_w=\tfrac12z_1^2+\tfrac12z_2^2+\tfrac{w}{2}\bigl(z_3^2+z_4^2\bigr)
+\text{(network terms)} .
\end{equation}
Substituting \eqref{eq:fz1}--\eqref{eq:fz4},
\[
\begin{aligned}
\dot z_1&=z_2+\chi_1-c_1z_1, &
\dot z_2&=(F_1-\hat F_1)+a_2z_3+a_2\chi_2-c_2z_2-z_1,\\
\dot z_3&=z_4+\chi_3-c_3z_3-\tfrac{a_2}{w}z_2, &
\dot z_4&=(F_2-\hat F_2)-c_4z_4-z_3,
\end{aligned}
\]
and the cross terms $z_1z_2$, $a_2z_2z_3$, and $w\,z_3z_4$ cancel pairwise in
$\dot{\mathcal{V}}_w$ exactly as in the proof of Theorem~\ref{thm:control}. The
network update laws \eqref{eq:upd2}--\eqref{eq:upd1} apply with $\varrho_1=z_2$
and $\varrho_2=w\,z_4$, the second network's learning rates being divided by $w$
so that its adaptation speed is independent of the weighting. The bound
\eqref{eq:Vdotfinal} follows with $A_2$ replaced by $a_2$ and the $z_3$, $z_4$
terms multiplied by $w$, under the same gain conditions \eqref{eq:gaincond}.

\rev{The skew coupling between $z_2$ and $z_3$ oscillates at
$\abs{a_2}/\sqrt{w}$. With $w=1$ that is about \SI{5.6}{rad/s}, inside the
command-filter bandwidths, and the linearized closed loop (networks off) is
stable at 672 of the 720 gain and bandwidth combinations examined in
Section~\ref{sec:cfinstab}, including the design gains; in case C2 the
nonlinear loop at the design gains balances with \SI{9.3}{mm} RMS position
error. With $w=a_2^2$ the coupling term becomes $z_2/a_2$, the coupling frequency
drops to about \SI{1}{rad/s}, and all 720 combinations are stable. That value is
used for the results of Chapter~\ref{ch:results} for its wider stability
margin over the gain grid, at some cost in tracking (\SI{16}{mm} RMS in C2).}

\subsection{What the coordinates change about the assumptions}

The filter-error assumption, Assumption~\ref{as:filt}, is now consistent with
the construction. The virtual controls depend on the state, the filter states,
and the network outputs, but not on the control, and the tracked output has no
zero dynamics. So bounded tracking errors imply bounded virtual controls and
bounded derivatives. \rev{It does not follow that the filter errors shrink as the
bandwidth is raised. The $O(1/\varpi_i)$ scaling of Assumption~\ref{as:filt}
holds for smooth virtual controls; with measurement noise the virtual controls
are themselves noisy, and a wider filter passes more of that noise.
Figure~\ref{fig:chi} shows both effects: $\chi_1$ and $\chi_2$ stay near
$10^{-3}$, while $\chi_3$, the error of the fastest filter, grows from
\num{0.005} to \num{0.18} over the bandwidths tested. The bandwidth therefore
trades filter lag against noise transmission, and $\bar\chi_i$ in the bound has
a floor set by the sensors.}

One modeling limitation remains. A parameter error in $k_m$ or $R$ changes the
true $B_2$, so the residual $F_2-\hat F_2$ contains $(B_2^{\mathrm{true}}-B_2)u$.
That residual depends on the control, and neither network is given the control
as an input. The approximation hypothesis therefore does not cover input-gain
error; robustness to it comes from the gains. \rev{It does not depend on a narrow
choice of them: in the closed-loop sweep of Section~\ref{sec:ctrlresults}, 63 of
64 gain sets balanced in both C2 and C4.}}
